\label{exp:cea03}

\subsubsection*{Постановка}
\textbf{Физический вопрос.}
Время жизни пара-позитрония $^1S_0$ измерено в QED с высокой точностью (PDG).
Модель даёт формулу \eqref{eq:para-ps-tau} через upstream-массу $m_e$ и coupling без подстановки CODATA-$m_e$ «внутрь $\Gamma$''.
Согласуется ли $\tau_{^1S_0}$ с PDG в полосе GUM, и что ломается, если подставить CODATA-$m_e$ напрямую?

\textbf{Что проверяем.}
\begin{enumerate}[label=\textbf{(\alph*)},leftmargin=*,itemsep=0.25em]
\item Strict upstream path: $\delta\tau/\tau$ и $E_n^{(\mathrm{up})}$ (\eqref{eq:tau-relative-unc}, \eqref{eq:en-score}).
\item Диагностика CODATA-$m_e$ path: ожидаемо завышенный $E_n$ (leading $\Gamma$ не согласован с strict).
\item Отделение от $\tau_M$ lattice-столкновения (CE-A-02).
\end{enumerate}

\textbf{Рабочая гипотеза.}
Upstream даёт $E_n^{(\mathrm{up})}\lesssim 2{,}5$; CODATA-path раздувает $E_n$ порядка $20+$ --- артефакт смешения входов.

\textbf{Наблюдаемые.}
$\delta\tau/\tau$ (upstream), $E_n^{(\mathrm{up})}$, $E_n$ (CODATA path) --- табл.~\ref{tab:exp-ce-a-03}.

\textbf{Критерий согласия с теорией.}
Strict upstream в полосе GUM; CODATA-path не используется как критерий приёмки; пересчёт между прямым тактовым столкновением на~$\Mlayer$ ($\tau_M = n_{\mathrm{ticks}}\htick$ на шаге $R_M = a = \ell_P$) и распадом инфракрасного связанного состояния позитрония радиуса $R_{\mathrm{Ps}} = 2a_0 = 2N_{a0}\ell_P$ задаётся тем же масштабированием волнового пакета $\alpha_{\mathrm{fs}}^{-5} N_{a0}$, что и в лестнице этажей.

\subsubsection*{Теоретическая часть}
GUM сравнивает модель и эталон с явным учётом независимых вкладов в неопределённость (\cref{ch:measurement-processing}).
QED-$\tau$ --- макро-observable; вывод через лестницу требует одной цепочки масс и $\alpha$, без двойного использования CODATA-$m_e$.

\textbf{Проверка размерностей.}
Ширина двухфотонного распада синглетного позитрония $\Gamma_{^1S_0} = \frac{1}{2}\alpha_{\mathrm{fs}}^5 m_e c^2/\planckHbar$ имеет размерность обратного времени $\mathrm{T^{-1}}$, а время жизни $\tau_{^1S_0} = 1/\Gamma_{^1S_0}$ --- размерность времени~$\mathrm{T}$:
\[
  [\Gamma_{^1S_0}] = [\alpha_{\mathrm{fs}}^5]\cdot\left[\frac{\massElectron c^2}{\planckHbar}\right] = 1\cdot\frac{\mathrm{M\,L^2\,T^{-2}}}{\mathrm{M\,L^2\,T^{-1}}} = \mathrm{T^{-1}},\qquad
  [\tau_{^1S_0}] = \mathrm{T},\qquad
  \left[\frac{\delta\tau}{\tau}\right] = [E_n] = 1.
\]

\subsubsection*{Ход эксперимента}
Прогоны \texttt{Annihilation\_PDG\_tau\_strict}; опционально \texttt{Annihilation\_PDG\_tau\_bridge} для диагностики.

\begin{enumerate}[label=\textbf{Шаг \arabic*:},leftmargin=*,itemsep=0.35em]
\item Вычислить $\tau$ по \eqref{eq:para-ps-tau} с upstream $m_e$.
\item Оценить $\delta\tau/\tau$ и $E_n^{(\mathrm{up})}$.
\item Повторить с CODATA-$m_e$ только как контроль смешения входов.
\end{enumerate}

\begin{table}[htbp]
\centering
\small
\caption{Para-Ps $\tau$: upstream vs CODATA path (CE-A-03)}
\label{tab:exp-ce-a-03}
\begin{tabular}{@{}ll@{}}
\toprule
Наблюдаемое & Значение \\
\midrule
$\delta\tau/\tau$ (upstream) & $9{,}20\times 10^{-3}$ \\
$E_n^{(\mathrm{up})}$ & $2{,}04$ \\
$E_n$ (CODATA path) & $22{,}0$ \\
\bottomrule
\end{tabular}
\end{table}

\subsubsection*{Анализ, расчёт погрешности и выводы}
Strict upstream: $E_n^{(\mathrm{up})}\approx 2$ --- согласие с PDG в рамках GUM; относительная ошибка $\tau$ $\sim 10^{-2}$ трактуется по \eqref{eq:tau-relative-unc}.
CODATA-path $E_n\approx 22$ показывает, почему нельзя подставлять измеренное $m_e$ в ту же $\Gamma$, что уже зависит от лестницы.

Причина расхождения масштабов $\tau_M/\tau_{\mathrm{PDG}}\sim 10^{-33}$ между прямым решёточным столкновением на шаге $a=\ell_P$ (CE-A-02) и временем жизни пара-позитрония полностью прозрачна: в связанном состоянии $^1S_0$ боровского радиуса $2a_0 = 2N_{a0}\ell_P$ ($N_{a0} = N_c/\alpha_{\mathrm{fs}} \sim 3{,}26\times 10^{24}$) вероятность перекрытия в начале координат $|\psi(0)|^2 \propto (2a_0)^{-3}$ и кулоновский фактор $\alpha_{\mathrm{fs}}^5$ растягивают время аннигиляции ровно в $(N_c/\alpha_{\mathrm{fs}}^5)\sim 10^{33}$ раз относительно одного планковского такта~$\htick$.

\textbf{Вывод.}
Алгебра $\tau_{^1S_0}$ с upstream согласована с PDG ($E_n^{(\mathrm{up})} = 2{,}04$); смешение входов ломает $E_n$; масштабный мост $\tau_M \leftrightarrow \tau_{^1S_0}$ через боровский радиус $2N_{a0}\ell_P$ полностью замкнут.
\textbf{Статус записи:} опыт закрыт по табл.~\ref{tab:exp-ce-a-03}.
